\documentclass[11pt,a4paper]{article}

% ── Encodage & langue ─────────────────────────────────────────────────────────
\usepackage[utf8]{inputenc}
\usepackage[T1]{fontenc}
\usepackage[french]{babel}

% ── Mise en page ──────────────────────────────────────────────────────────────
\usepackage[top=2.2cm, bottom=2.2cm, left=2.5cm, right=2.5cm]{geometry}
\usepackage{parskip}
\usepackage{enumitem}
\usepackage{booktabs}
\usepackage{array}

% ── Typo ──────────────────────────────────────────────────────────────────────
\usepackage{lmodern}
\usepackage{microtype}

% ── Couleurs & boîtes ─────────────────────────────────────────────────────────
\usepackage[dvipsnames]{xcolor}
\usepackage{tcolorbox}
\tcbuselibrary{skins, breakable, listings}

\definecolor{bleu}{RGB}{30, 80, 160}
\definecolor{vertTP}{RGB}{0, 120, 80}
\definecolor{orangeHint}{RGB}{200, 100, 0}
\definecolor{griscode}{RGB}{248, 248, 248}
\definecolor{rougeErr}{RGB}{180, 30, 30}

\newtcolorbox{objectif}{
  colback=bleu!8, colframe=bleu, fonttitle=\bfseries,
  title=Objectif, breakable
}
\newtcolorbox{exercice}[1]{
  colback=vertTP!6, colframe=vertTP, fonttitle=\bfseries,
  title={Exercice #1}, breakable
}
\newtcolorbox{hint}{
  colback=orangeHint!8, colframe=orangeHint, fonttitle=\bfseries\small,
  title={Aide}, breakable
}
\newtcolorbox{validation}{
  colback=gray!8, colframe=gray!60, fonttitle=\bfseries,
  title={Validation}, breakable
}
\newtcolorbox{attention}{
  colback=rougeErr!8, colframe=rougeErr, fonttitle=\bfseries,
  title={Attention}, breakable
}

% ── Code ──────────────────────────────────────────────────────────────────────
\usepackage{listings}

\lstdefinestyle{python}{
  language=Python,
  backgroundcolor=\color{griscode},
  basicstyle=\ttfamily\footnotesize,
  keywordstyle=\color{bleu}\bfseries,
  stringstyle=\color{vertTP},
  commentstyle=\color{gray!70}\itshape,
  numberstyle=\tiny\color{gray},
  numbers=left,
  numbersep=8pt,
  breaklines=true,
  showstringspaces=false,
  frame=single,
  framerule=0.4pt,
  rulecolor=\color{gray!40},
  tabsize=4,
  morekeywords={as, with, from, import, None, True, False},
  literate={é}{{\'{e}}}1 {è}{{\`{e}}}1 {ê}{{\^{e}}}1 {à}{{\`{a}}}1
           {â}{{\^{a}}}1 {ù}{{\`{u}}}1 {î}{{\^{i}}}1 {ô}{{\^{o}}}1
           {ç}{{\c{c}}}1 {→}{{$\rightarrow$}}1
}
\lstdefinestyle{bash}{
  language=bash,
  backgroundcolor=\color{griscode},
  basicstyle=\ttfamily\footnotesize\color{black},
  keywordstyle=\color{bleu},
  commentstyle=\color{gray!70}\itshape,
  breaklines=true,
  frame=single,
  framerule=0.4pt,
  rulecolor=\color{gray!40},
  numbers=none,
  showstringspaces=false,
}
\lstset{style=python}

% ── Titres ────────────────────────────────────────────────────────────────────
\usepackage{titlesec}
\titleformat{\section}{\large\bfseries\color{bleu}}{}{0em}{\thesection\quad}[\titlerule]
\titleformat{\subsection}{\normalsize\bfseries\color{bleu!80}}{}{0em}{\thesubsection\quad}

% ── Diagrammes ────────────────────────────────────────────────────────────────
\usepackage{tikz}
\usetikzlibrary{shapes.geometric, arrows.meta, positioning, fit, backgrounds}

% ── En-tête / pied ────────────────────────────────────────────────────────────
\usepackage{fancyhdr}
\pagestyle{fancy}
\fancyhf{}
\lhead{\small\textcolor{bleu}{\textbf{Pipeline ETL — Répertoire National des Élus}}}
\rhead{\small\textcolor{gray}{BAC5 — Data Engineering}}
\cfoot{\small\thepage}
\renewcommand{\headrulewidth}{0.4pt}

\usepackage{hyperref}
\hypersetup{colorlinks=true, linkcolor=bleu, urlcolor=bleu}

% =============================================================================
\begin{document}

% ── Page de titre ─────────────────────────────────────────────────────────────
\begin{titlepage}
  \centering
  \vspace*{2cm}
  {\LARGE\bfseries\color{bleu} Travaux Pratiques\\[0.4em]
   Pipeline ETL — Répertoire National des Élus\par}
  \vspace{0.8cm}
  {\large\color{gray} Apache Spark \textbullet{} HDFS \textbullet{} Apache Airflow \textbullet{} PostgreSQL\par}
  \vspace{2.5cm}
\begin{center}
  
  \begin{tikzpicture}[scale=0.9, transform shape,
    node distance=1.6cm,
    box/.style={rectangle, rounded corners=4pt, draw, minimum width=2.6cm,
                minimum height=1cm, align=center, font=\small\bfseries},
    arr/.style={-Stealth, thick, color=bleu!60},
  ]
    \node[box, fill=bleu!12, draw=bleu]    (pg)  {PostgreSQL\\{\tiny 12 tables RNE}};
    \node[box, fill=orange!15, draw=orange, right=of pg]  (raw) {HDFS\\{\tiny /rne/raw}};
    \node[box, fill=orange!22, draw=orange!80, right=of raw] (ref) {HDFS\\{\tiny /rne/refined}};
    \node[box, fill=orange!30, draw=orange!60, right=of ref] (gld) {HDFS\\{\tiny /rne/gold}};
    \node[box, fill=bleu!12, draw=bleu, right=of gld]  (pgout) {PostgreSQL\\{\tiny 4 agrégations}};

    \draw[arr] (pg)  -- node[above,font=\tiny\bfseries]{Extract}   (raw);
    \draw[arr] (raw) -- node[above,font=\tiny\bfseries]{Transform} (ref);
    \draw[arr] (ref) -- node[above,font=\tiny\bfseries]{Aggregate} (gld);
    \draw[arr] (gld) -- node[above,font=\tiny\bfseries]{Write}     (pgout);

    \node[draw=vertTP, dashed, rounded corners, fit=(pg)(pgout),
          inner sep=12pt,
          label={[font=\small\bfseries\color{vertTP}]above:Airflow DAG — orchestration}] {};
  \end{tikzpicture}
\end{center}

  \vspace{2.5cm}
  \begin{tabular}{ll}
    \toprule
    \textbf{Durée}   & 3h \\
    \textbf{Données} & Répertoire National des Élus — data.gouv.fr \\
    \textbf{Stack}   & Python 3.12 · PySpark 3.5 · Airflow 2.x · PostgreSQL 15 \\
    \bottomrule
  \end{tabular}
\end{titlepage}

% \tableofcontents
\newpage

% =============================================================================
\section{Contexte et objectifs}
% =============================================================================

Le \textbf{Répertoire National des Élus (RNE)} est un jeu de données publié par le Ministère de l'Intérieur sur \texttt{data.gouv.fr}. Il recense l'ensemble des mandats électoraux en France : maires, conseillers municipaux, députés, sénateurs, etc.

Vous allez construire un \textbf{pipeline ETL complet} en 4 étapes :
\begin{enumerate}
  \item \textbf{Extract} — lire les 12 tables PostgreSQL et les écrire en Parquet sur HDFS ;
  \item \textbf{Transform} — unifier les 12 tables en un seul DataFrame \texttt{elus\_unified} ;
  \item \textbf{Aggregate} — calculer 4 agrégations analytiques et les écrire dans PostgreSQL ;
  \item \textbf{Orchestrate} — assembler le tout dans un DAG Airflow.
\end{enumerate}

\subsection*{Questions analytiques}

\begin{center}
\begin{tabular}{clll}
  \toprule
  & \textbf{Question} & \textbf{Table} & \textbf{Graphique} \\
  \midrule
  Q1 & Répartition H/F par type de mandat       & \texttt{agg\_parite}     & Barres empilées \\
  Q2 & CSP dominantes par type de mandat         & \texttt{agg\_csp}        & Treemap \\
  Q3 & Durée moyenne de mandat par région        & \texttt{agg\_anciennete} & Carte choroplèthe \\
  Q5 & Distribution d'âge par type de mandat     & \texttt{agg\_age}        & Boîte à moustaches \\
  \bottomrule
\end{tabular}
\end{center}

% =============================================================================
\section{Structure du projet}
% =============================================================================

Fichiers fournis (\textcolor{gray}{gris}) et fichiers à créer (\textcolor{vertTP}{\textbf{vert}}) :

\begin{lstlisting}[style=bash, numbers=none]
.
├── dags/
│   └── rne_pipeline.py            # (à créer) DAG Airflow
├── infra/
│   ├── docker-compose.yml         # fourni
│   ├── init-postgres.sql          # fourni — DDL des 12 tables
│   └── migrate_csv_to_postgres.py # fourni — chargement des CSV
├── rep_national_elus/
│   └── *.csv                      # 12 fichiers CSV
└── spark_jobs/
    ├── __init__.py                # fourni (vide)
    ├── config.py                  # fourni — variables + build_spark()
    ├── job_extract.py             # (à créer)
    ├── job_transform.py           # (à créer)
    └── job_aggregate.py           # (à créer)
\end{lstlisting}

\subsection*{Ce que contient \texttt{config.py}}

Tout ce dont vous avez besoin est déjà défini dans \texttt{config.py} — vous n'avez qu'à importer :

\begin{lstlisting}
# Connexion PostgreSQL
JDBC_URL    # url jdbc:postgresql://...
JDBC_PROPS  # {"user": ..., "password": ..., "driver": ...}

# Chemins HDFS
HDFS_RAW      # hdfs://namenode:9000/rne/raw
HDFS_REFINED  # hdfs://namenode:9000/rne/refined
HDFS_GOLD     # hdfs://namenode:9000/rne/gold

# Listes de tables
TABLES_SOURCES  # liste des 12 noms de tables
TABLES_AGG      # ["agg_parite", "agg_csp", "agg_anciennete", "agg_age"]
TABLE_SPECS     # mapping géographique (voir Exercice 2)

# Fonction utilitaire
build_spark(app_name)  # crée et retourne une SparkSession configurée
\end{lstlisting}

% =============================================================================
\section{Démarrage de l'infrastructure}
% =============================================================================

\begin{objectif}
Lancer la stack Docker et vérifier que PostgreSQL est chargé avant d'écrire le moindre code.
\end{objectif}

Depuis le dossier \texttt{infra/} :
\begin{lstlisting}[style=bash, numbers=none]
docker compose up -d
\end{lstlisting}

\begin{attention}
Le service \texttt{init-db} charge automatiquement les 12 CSV dans PostgreSQL au premier démarrage. Cette opération prend \textbf{5 à 10 minutes} (environ 485 000 lignes). Attendez que le service soit \texttt{exited (0)} avant de continuer.
\end{attention}

Vérifiez :
\begin{lstlisting}[style=bash, numbers=none]
docker compose logs init-db | tail -5
\end{lstlisting}

\subsection*{Interfaces disponibles}

\begin{center}
\begin{tabular}{lll}
  \toprule
  \textbf{Service}  & \textbf{URL}                       & \textbf{Accès} \\
  \midrule
  Airflow UI        & \texttt{http://localhost:8080}      & admin / admin \\
  Spark Master UI   & \texttt{http://localhost:8090}      & — \\
  HDFS NameNode UI  & \texttt{http://localhost:9870}      & — \\
  \bottomrule
\end{tabular}
\end{center}

% =============================================================================
\section{Exercice 1 — Extraction PostgreSQL \texorpdfstring{$\rightarrow$}{→} HDFS}
% =============================================================================

\textbf{Fichier :} \texttt{spark\_jobs/job\_extract.py}

\begin{objectif}
Lire les 12 tables du schéma \texttt{rne} depuis PostgreSQL via JDBC et écrire chaque table en Parquet sur HDFS dans la zone \texttt{raw}.
\end{objectif}

\begin{lstlisting}
import logging

from config import HDFS_RAW, JDBC_PROPS, JDBC_URL, TABLES_SOURCES, build_spark

logging.basicConfig(level=logging.INFO, format="%(asctime)s [extract] %(message)s")
logger = logging.getLogger(__name__)


def run_extract():
    spark = build_spark("ETL-RNE-Extract")
    spark.sparkContext.setLogLevel("WARN")
    try:
        for table in TABLES_SOURCES:
            # TODO 1 : Lire la table rne.<table> depuis PostgreSQL via JDBC
            #   Utiliser : spark.read.jdbc(url=..., table=..., properties=...)
            df = ...

            # TODO 2 : Ecrire df en Parquet (mode overwrite) vers HDFS_RAW/<table>
            ...

            logger.info(f"[extract] {table} -> {df.count():,} lignes")
        logger.info("[extract] Terminé.")
    finally:
        spark.stop()


if __name__ == "__main__":
    run_extract()
\end{lstlisting}

\begin{hint}
\texttt{spark.read.jdbc(url=JDBC\_URL, table="rne.nom\_table", properties=JDBC\_PROPS)} lit une table PostgreSQL.\\
\texttt{df.write.mode("overwrite").parquet(chemin)} écrit en Parquet.
\end{hint}

\begin{validation}
Après exécution, l'interface HDFS (\texttt{http://localhost:9870}) doit afficher 12 dossiers dans \texttt{/rne/raw/}.
\end{validation}

% =============================================================================
\section{Exercice 2 — Transformation et unification}
% =============================================================================

\textbf{Fichier :} \texttt{spark\_jobs/job\_transform.py}

\begin{objectif}
Lire les 12 tables Parquet depuis la zone \texttt{raw}, les unifier en un seul DataFrame \texttt{elus\_unified} avec un schéma commun, et l'écrire dans la zone \texttt{refined}.
\end{objectif}

\subsection*{Schéma cible}

\begin{center}
\begin{tabular}{lll}
  \toprule
  \textbf{Colonne} & \textbf{Type} & \textbf{Remarque} \\
  \midrule
  \texttt{type\_mandat}         & string  & ex : "Maire", "Député" \\
  \texttt{nom}, \texttt{prenom} & string  & \\
  \texttt{code\_sexe}           & string  & "M" ou "F" \\
  \texttt{date\_naissance}      & date    & \\
  \texttt{age}                  & integer & calculé à aujourd'hui \\
  \texttt{code\_csp}            & integer & \\
  \texttt{libelle\_csp}         & string  & \\
  \texttt{date\_debut\_mandat}  & date    & \\
  \texttt{duree\_mandat\_jours} & integer & calculé à aujourd'hui \\
  \texttt{code\_departement}    & string  & NULL si absent \\
  \texttt{libelle\_departement} & string  & NULL si absent \\
  \texttt{code\_region}         & string  & NULL si absent \\
  \texttt{libelle\_region}      & string  & NULL si absent \\
  \bottomrule
\end{tabular}
\end{center}

\subsection*{Pourquoi \texttt{TABLE\_SPECS} ?}

Les 12 tables n'ont pas les mêmes colonnes géographiques. Par exemple, la table \texttt{conseillers\_regionaux} appelle la colonne département \texttt{code\_section\_departementale}, tandis que les maires utilisent \texttt{code\_departement}. \texttt{TABLE\_SPECS} (fourni dans \texttt{config.py}) décrit ce mapping pour chaque table :

\begin{lstlisting}
# Extrait de config.py — fourni, ne pas réécrire
TABLE_SPECS = [
    ("conseillers_municipaux", "Conseiller municipal",
     {"code_region": None,          # absent -> NULL
      "libelle_region": None,
      "code_departement": "code_departement",
      "libelle_departement": "libelle_departement"}),

    ("conseillers_regionaux", "Conseiller régional",
     {"code_region": "code_region",
      "libelle_region": "libelle_region",
      "code_departement": "code_section_departementale",  # nom different
      "libelle_departement": "libelle_section_departementale"}),
    # ... 10 autres tables
]
\end{lstlisting}

\subsection*{Squelette à compléter}

\begin{lstlisting}
import logging

from pyspark.sql import functions as F
from config import HDFS_RAW, HDFS_REFINED, TABLE_SPECS, build_spark

logging.basicConfig(level=logging.INFO, format="%(asctime)s [transform] %(message)s")
logger = logging.getLogger(__name__)


def build_unified_df(spark):
    TODAY = F.current_date()  # IMPORTANT : dans la fonction, pas au niveau module
    dfs = []

    for table_name, mandat_label, geo in TABLE_SPECS:
        df = spark.read.parquet(f"{HDFS_RAW}/{table_name}")

        # TODO 1 : Construire la liste `select` avec les colonnes communes :
        #   F.lit(mandat_label).alias("type_mandat")
        #   F.col("nom"), F.col("prenom"), F.col("code_sexe")
        #   F.col("date_naissance")
        #   age  = datediff(TODAY, date_naissance) / 365.25  cast("integer")
        #   F.col("code_csp").cast("integer"), F.col("libelle_csp")
        #   F.col("date_debut_mandat")
        #   duree_mandat_jours = datediff(TODAY, date_debut_mandat)
        select = [...]

        # TODO 2 : Pour chaque (col_cible, col_source) dans geo.items() :
        #   Si col_source existe dans df.columns -> F.col(col_source).alias(col_cible)
        #   Sinon                                -> F.lit(None).cast("string").alias(col_cible)
        for col_cible, col_source in geo.items():
            ...

        dfs.append(df.select(select))

    # TODO 3 : Unifier tous les DataFrames avec unionByName()
    unified = dfs[0]
    for df in dfs[1:]:
        unified = ...

    # TODO 4 : Filtrer les lignes où nom ou date_naissance est NULL
    return unified.filter(...)


def run_transform():
    spark = build_spark("ETL-RNE-Transform")
    spark.sparkContext.setLogLevel("WARN")
    try:
        unified = build_unified_df(spark)
        logger.info(f"[transform] {unified.count():,} lignes unifiées")
        unified.write.mode("overwrite").parquet(f"{HDFS_REFINED}/elus_unified")
        logger.info("[transform] Terminé.")
    finally:
        spark.stop()


if __name__ == "__main__":
    run_transform()
\end{lstlisting}

\begin{hint}
\begin{itemize}[leftmargin=*]
  \item \texttt{F.datediff(d1, d2)} retourne le nombre de jours entre deux dates.
  \item \texttt{F.lit(None).cast("string")} crée une colonne NULL de type string.
  \item \texttt{unionByName()} fusionne deux DataFrames en alignant par \textbf{nom} de colonne.
  \item \texttt{F.col("nom").isNotNull()} teste si une colonne est non-NULL.
\end{itemize}
\end{hint}

\begin{attention}
\texttt{F.current\_date()} nécessite une SparkSession active. Ne l'appelez \textbf{jamais au niveau du module} (hors de toute fonction) — cela provoque une \texttt{AssertionError} au démarrage du script.
\end{attention}

% =============================================================================
\section{Exercice 3 — Agrégations}
% =============================================================================

\textbf{Fichier :} \texttt{spark\_jobs/job\_aggregate.py}

\begin{objectif}
Lire \texttt{elus\_unified} depuis la zone \texttt{refined}, calculer les 4 agrégations, et les écrire dans la zone \texttt{gold} (Parquet) \textbf{et} dans PostgreSQL.
\end{objectif}

\begin{lstlisting}
import logging

from pyspark.sql import functions as F
from pyspark.sql.window import Window
from config import HDFS_GOLD, HDFS_REFINED, JDBC_PROPS, JDBC_URL, build_spark

logging.basicConfig(level=logging.INFO, format="%(asctime)s [aggregate] %(message)s")
logger = logging.getLogger(__name__)


def write_to_postgres(df, table_name):
    """Ecrit df sur HDFS (gold) et dans PostgreSQL."""
    df.write.mode("overwrite").parquet(f"{HDFS_GOLD}/{table_name}")
    df.write.mode("overwrite").jdbc(
        url=JDBC_URL, table=f"rne.{table_name}", properties=JDBC_PROPS
    )
    logger.info(f"[aggregate] {table_name} -> {df.count():,} lignes")


# Q1 — Répartition H/F par type de mandat
# Résultat : type_mandat | code_sexe | nb_elus | pct
# pct = % de ce sexe dans le type_mandat (utiliser Window.partitionBy)
def agg_parite(df):
    # TODO
    ...


# Q2 — CSP dominantes par type de mandat
# Résultat : type_mandat | code_csp | libelle_csp | nb_elus
# Filtrer libelle_csp NULL. Trier par nb_elus décroissant.
def agg_csp(df):
    # TODO
    ...


# Q3 — Durée moyenne de mandat par région
# Résultat : type_mandat | code_region | libelle_region
#          | duree_moyenne_jours | duree_moyenne_ans | nb_elus
# duree_moyenne_ans = duree_moyenne_jours / 365.25
def agg_anciennete(df):
    # TODO
    ...


# Q5 — Distribution d'âge par type de mandat
# Résultat : type_mandat | age_moyen | age_min | age_max
#          | q1_age | mediane_age | q3_age | nb_elus
# Utiliser F.percentile_approx("age", 0.25 / 0.50 / 0.75)
def agg_age(df):
    # TODO
    ...


def run_aggregate():
    spark = build_spark("ETL-RNE-Aggregate")
    spark.sparkContext.setLogLevel("WARN")
    try:
        df = spark.read.parquet(f"{HDFS_REFINED}/elus_unified")
        df.cache()
        logger.info(f"[aggregate] {df.count():,} lignes lues")

        write_to_postgres(agg_parite(df),     "agg_parite")
        write_to_postgres(agg_csp(df),        "agg_csp")
        write_to_postgres(agg_anciennete(df), "agg_anciennete")
        write_to_postgres(agg_age(df),        "agg_age")

        logger.info("[aggregate] Terminé.")
    finally:
        spark.stop()


if __name__ == "__main__":
    run_aggregate()
\end{lstlisting}

\begin{hint}
\begin{itemize}[leftmargin=*]
  \item \texttt{Window.partitionBy("col")} calcule un total par groupe sans \texttt{groupBy}.
  \item \texttt{F.percentile\_approx("age", 0.5)} retourne la médiane (efficace sur grand volume).
  \item \texttt{F.round(expr, 2)} arrondit à 2 décimales.
  \item \texttt{df.cache()} garde le DataFrame en mémoire pour éviter de le relire 4 fois.
\end{itemize}
\end{hint}

% =============================================================================
\section{Exercice 4 — DAG Airflow}
\label{sec:dag}
% =============================================================================

\textbf{Fichier :} \texttt{dags/rne\_pipeline.py}

\begin{objectif}
Créer un DAG Airflow qui orchestre le pipeline en 5 tâches séquentielles.
\end{objectif}

\subsection*{Architecture}

\begin{center}
\begin{tikzpicture}[
  node distance=1.2cm,
  task/.style={rectangle, rounded corners=3pt, draw=bleu, fill=bleu!10,
               minimum width=3cm, minimum height=0.8cm,
               align=center, font=\small\bfseries\color{bleu}},
  arr/.style={-Stealth, thick, color=bleu!50}
]
  \node[task] (t1) {check\_source\\{\tiny PythonOperator}};
  \node[task, right=of t1] (t2) {job\_extract\\{\tiny SparkSubmit}};
  \node[task, right=of t2] (t3) {job\_transform\\{\tiny SparkSubmit}};
  \node[task, right=of t3] (t4) {job\_aggregate\\{\tiny SparkSubmit}};
  \node[task, right=of t4] (t5) {validate\_output\\{\tiny PythonOperator}};
  \draw[arr](t1)--(t2); \draw[arr](t2)--(t3);
  \draw[arr](t3)--(t4); \draw[arr](t4)--(t5);
\end{tikzpicture}
\end{center}

\begin{lstlisting}
import logging
from datetime import datetime, timedelta

from airflow import DAG
from airflow.operators.python import PythonOperator
from airflow.providers.apache.spark.operators.spark_submit import SparkSubmitOperator

logger = logging.getLogger(__name__)

default_args = {
    "owner": "bac5",
    "retries": 1,
    "retry_delay": timedelta(minutes=2),
    "email_on_failure": False,
}

SPARK_CONF = {"spark.executor.memory": "2g", "spark.driver.memory": "1g"}

# TODO 1 : Définir SPARK_KWARGS (dict) utilisé par les 3 SparkSubmitOperator
#   conn_id  = "spark_default"
#   jars     = "/opt/airflow/jars/postgresql-42.7.3.jar"
#   conf     = SPARK_CONF
#   py_files = les deux fichiers Python à embarquer dans chaque job Spark :
#              __init__.py et config.py (chemin complet dans /opt/airflow/spark_jobs/)
SPARK_KWARGS = dict(...)


def check_source(**ctx):
    """Vérifie que les 12 tables sources sont non vides."""
    import psycopg2
    from spark_jobs.config import (
        POSTGRES_DB, POSTGRES_HOST, POSTGRES_PASS,
        POSTGRES_PORT, POSTGRES_USER, TABLES_SOURCES,
    )
    conn = psycopg2.connect(
        host=POSTGRES_HOST, port=POSTGRES_PORT,
        user=POSTGRES_USER, password=POSTGRES_PASS, dbname=POSTGRES_DB,
    )
    cur = conn.cursor()
    for table in TABLES_SOURCES:
        cur.execute(f"SELECT COUNT(*) FROM rne.{table}")
        count = cur.fetchone()[0]
        if count == 0:
            raise ValueError(f"Table rne.{table} vide !")
        logger.info(f"[check] rne.{table} : {count:,} lignes")
    conn.close()


def validate_output(**ctx):
    """Vérifie que les 4 tables d'agrégation sont peuplées."""
    import psycopg2
    from spark_jobs.config import (
        POSTGRES_DB, POSTGRES_HOST, POSTGRES_PASS,
        POSTGRES_PORT, POSTGRES_USER, TABLES_AGG,
    )
    conn = psycopg2.connect(
        host=POSTGRES_HOST, port=POSTGRES_PORT,
        user=POSTGRES_USER, password=POSTGRES_PASS, dbname=POSTGRES_DB,
    )
    cur = conn.cursor()
    for table in TABLES_AGG:
        cur.execute(f"SELECT COUNT(*) FROM rne.{table}")
        count = cur.fetchone()[0]
        if count == 0:
            raise ValueError(f"Table rne.{table} vide !")
        logger.info(f"[validate] rne.{table} : {count:,} lignes")
    conn.close()
    logger.info("[validate] Pipeline validé.")


# TODO 2 : Définir le DAG
#   dag_id           = "ETL-DAG-RNE"
#   start_date       = datetime(2026, 1, 1)
#   catchup          = False
#   schedule_interval = None   (déclenchement manuel uniquement)
#   default_args     = default_args
#   tags             = ["ETL", "RNE", "Spark"]
with DAG(...) as dag:

    # TODO 3 : Créer les 5 tâches
    #
    #   health_check  = PythonOperator(task_id="check_source",    ...)
    #   extract_job   = SparkSubmitOperator(task_id="job_extract",
    #                     application="/opt/airflow/spark_jobs/job_extract.py",
    #                     name="ETL-RNE-Extract",
    #                     execution_timeout=timedelta(hours=1),
    #                     **SPARK_KWARGS)
    #   transform_job = SparkSubmitOperator(...)  # job_transform.py
    #   aggregate_job = SparkSubmitOperator(...)  # job_aggregate.py
    #   verify_output = PythonOperator(task_id="validate_output", ...)

    health_check  = ...
    extract_job   = ...
    transform_job = ...
    aggregate_job = ...
    verify_output = ...

# TODO 4 : Définir les dépendances séquentielles
... >> ... >> ... >> ... >> ...
\end{lstlisting}

\begin{hint}
\begin{itemize}[leftmargin=*]
  \item \texttt{py\_files} permet d'envoyer des fichiers Python supplémentaires au job Spark.\\
        Format : \texttt{"chemin1.py,chemin2.py"} (séparés par une virgule, sans espace).
  \item \texttt{schedule\_interval=None} : le DAG ne se lance \textbf{jamais automatiquement}.\\
        Déclenchez-le manuellement via le bouton \textbf{▶ Trigger DAG} dans Airflow.
  \item \texttt{tags} appartient au constructeur \texttt{DAG(...)}, pas à \texttt{default\_args}.
\end{itemize}
\end{hint}

% =============================================================================
\section{Lancement et validation}
% =============================================================================

\subsection*{Déclencher le DAG}

Dans Airflow (\texttt{http://localhost:8080}), naviguez vers \texttt{ETL-DAG-RNE} et cliquez sur \textbf{▶ Trigger DAG}, ou depuis le terminal :

\begin{lstlisting}[style=bash, numbers=none]
docker compose exec airflow-scheduler airflow dags trigger ETL-DAG-RNE
\end{lstlisting}

\subsection*{Durées typiques}

\begin{center}
\begin{tabular}{lll}
  \toprule
  \textbf{Tâche}               & \textbf{Durée} & \textbf{Ce qu'elle fait} \\
  \midrule
  \texttt{check\_source}    & $<$ 5 s      & Compte les lignes dans les 12 tables PG \\
  \texttt{job\_extract}     & 1 -- 3 min   & Extrait 12 tables PG $\to$ Parquet HDFS \\
  \texttt{job\_transform}   & 1 -- 2 min   & Unifie en \texttt{elus\_unified} \\
  \texttt{job\_aggregate}   & 1 -- 2 min   & Calcule et écrit les 4 agrégations \\
  \texttt{validate\_output} & $<$ 5 s      & Vérifie les 4 tables résultat dans PG \\
  \bottomrule
\end{tabular}
\end{center}

\subsection*{Vérification finale}

Si toutes les tâches sont vertes, vérifiez les résultats dans PostgreSQL :

\begin{lstlisting}[style=bash, numbers=none]
-- Parité H/F chez les maires
SELECT code_sexe, nb_elus, pct FROM rne.agg_parite
WHERE type_mandat = 'Maire' ORDER BY code_sexe;

-- Top 5 CSP chez les députés
SELECT libelle_csp, nb_elus FROM rne.agg_csp
WHERE type_mandat = 'Député' ORDER BY nb_elus DESC LIMIT 5;
\end{lstlisting}

\subsection*{En cas d'erreur}

Cliquez sur la tâche rouge dans Airflow $\to$ onglet \textbf{Log}, ou :
\begin{lstlisting}[style=bash, numbers=none]
docker compose logs airflow-scheduler | grep -i "error\|failed"
\end{lstlisting}

% =============================================================================
\section{Exercice 5 — Visualisation avec Tableau Desktop}
% =============================================================================

\begin{objectif}
Connecter Tableau Desktop à PostgreSQL, créer 4 visualisations répondant aux questions analytiques, puis assembler un Dashboard interactif.
\end{objectif}

% ── 5.1 Connexion PostgreSQL ─────────────────────────────────────────────────
\subsection*{5.1 — Connexion à PostgreSQL}

\begin{enumerate}[leftmargin=*]
  \item Ouvrez \textbf{Tableau Desktop} et cliquez sur \textbf{Connecter $\to$ À un serveur $\to$ PostgreSQL}.
  \item Renseignez les paramètres suivants :
\end{enumerate}

\begin{center}
\begin{tabular}{ll}
  \toprule
  \textbf{Champ}     & \textbf{Valeur} \\
  \midrule
  Serveur            & \texttt{localhost} \\
  Port               & \texttt{5432} \\
  Base de données    & \texttt{rne} \\
  Schéma             & \texttt{rne} \\
  Nom d'utilisateur  & \texttt{spark} \\
  Mot de passe       & \texttt{spark} \\
  \bottomrule
\end{tabular}
\end{center}

\begin{hint}
Le conteneur PostgreSQL expose le port \texttt{5432} sur \texttt{localhost}. Vérifiez qu'il est bien démarré (\texttt{docker compose ps}) avant d'ouvrir Tableau.
\end{hint}

\begin{attention}
Les 4 tables d'agrégation doivent être présentes dans PostgreSQL \textbf{avant} de commencer cet exercice. Assurez-vous que le DAG Airflow (Exercice 4) s'est terminé avec succès.
\end{attention}

Une fois connecté, faites glisser les 4 tables dans l'espace de travail :
\texttt{rne.agg\_parite}, \texttt{rne.agg\_csp}, \texttt{rne.agg\_anciennete}, \texttt{rne.agg\_age}.

% ── 5.2 Feuilles de visualisation ───────────────────────────────────────────
\subsection*{5.2 — Création des feuilles de visualisation}

Créez \textbf{4 feuilles} distinctes, une par question analytique.

\begin{exercice}{5a — Q1 : Répartition H/F par type de mandat (Barres empilées)}
\textbf{Source :} \texttt{agg\_parite} — colonnes : \texttt{type\_mandat}, \texttt{code\_sexe}, \texttt{nb\_elus}, \texttt{pct}

\begin{enumerate}[leftmargin=*]
  \item Glissez \texttt{type\_mandat} dans \textbf{Colonnes}.
  \item Glissez \texttt{nb\_elus} dans \textbf{Lignes}.
  \item Glissez \texttt{code\_sexe} dans \textbf{Couleur} (volet \textit{Repères}).
  \item Dans le menu \textbf{Repères}, sélectionnez le type \textbf{Barre}.
  \item Cliquez sur \textbf{Analyse $\to$ Empiler les repères $\to$ On}.
  \item Ajoutez \texttt{pct} en \textbf{Étiquette} pour afficher le pourcentage sur chaque barre.
  \item Renommez la feuille \textit{``Parité H/F par mandat''}.
\end{enumerate}

\begin{validation}
Le graphique doit montrer pour chaque type de mandat deux barres empilées (M et F) dont la hauteur totale correspond au nombre total d'élus. Les pourcentages sont visibles sur chaque segment.
\end{validation}
\end{exercice}

\begin{exercice}{5b — Q2 : CSP dominantes par type de mandat (Treemap)}
\textbf{Source :} \texttt{agg\_csp} — colonnes : \texttt{type\_mandat}, \texttt{code\_csp}, \texttt{libelle\_csp}, \texttt{nb\_elus}

\begin{hint}
Les lignes avec \texttt{libelle\_csp} NULL sont déjà filtrées par le job Spark — la table ne les contient pas.
\end{hint}

\begin{enumerate}[leftmargin=*]
  \item Glissez \texttt{nb\_elus} dans \textbf{Taille} et dans \textbf{Couleur}.
  \item Glissez \texttt{type\_mandat} dans \textbf{Détail}.
  \item Glissez \texttt{libelle\_csp} dans \textbf{Étiquette}.
  \item Dans le menu \textbf{Repères}, sélectionnez \textbf{Carré}.
  \item Choisissez une palette séquentielle (ex. \textit{Bleu-Vert}) via \textbf{Modifier les couleurs}.
  \item Ajoutez \texttt{type\_mandat} dans \textbf{Filtres} pour explorer mandat par mandat.
  \item Renommez la feuille \textit{``CSP par mandat''}.
\end{enumerate}

\begin{validation}
Chaque rectangle représente une CSP ; les plus grands correspondent aux CSP les plus représentées. Le graphique est lisible sans valeurs NULL.
\end{validation}
\end{exercice}

\begin{exercice}{5c — Q3 : Durée moyenne de mandat par région (Carte choroplèthe)}
\textbf{Source :} \texttt{agg\_anciennete}

\begin{attention}
La table \texttt{agg\_anciennete} ne contient que les mandats pour lesquels \texttt{libelle\_region} est non NULL. Dans les données RNE, seuls les \textbf{Conseillers régionaux} et les \textbf{Membres d'assemblée} ont une colonne région renseignée — c'est donc la carte pour ces deux types de mandats uniquement.
\end{attention}

\begin{enumerate}[leftmargin=*]
  \item Double-cliquez sur \texttt{libelle\_region} — Tableau génère automatiquement une carte de France.
  \item Glissez \texttt{duree\_moyenne\_ans} dans \textbf{Couleur}.
  \item Glissez \texttt{nb\_elus} dans \textbf{Taille} (bulles proportionnelles) ou laissez le type \textbf{Rempli} pour une carte choroplèthe pure.
  \item Ajoutez \texttt{libelle\_region}, \texttt{duree\_moyenne\_ans} et \texttt{duree\_moyenne\_jours} dans \textbf{Info-bulle}.
  \item Ajoutez \texttt{type\_mandat} dans \textbf{Filtres} et activez \textbf{Afficher le filtre}.
  \item Renommez la feuille \textit{``Ancienneté par région''}.
\end{enumerate}

\begin{hint}
Si Tableau ne reconnaît pas les noms de régions, assignez manuellement le rôle géographique : clic droit sur \texttt{libelle\_region} $\to$ \textbf{Rôle géographique $\to$ État/Province}. Vous pouvez aussi utiliser \texttt{code\_region} (ex. \texttt{11} pour Île-de-France) avec le rôle \textbf{Code d'État/province}.
\end{hint}

\begin{validation}
La carte de France est colorée selon \texttt{duree\_moyenne\_ans}. Seules les régions liées aux mandats régionaux ont des données. Un filtre par \texttt{type\_mandat} est accessible.
\end{validation}
\end{exercice}

\begin{exercice}{5d — Q5 : Distribution d'âge par type de mandat (Boîte à moustaches)}
\textbf{Source :} \texttt{agg\_age} — colonnes : \texttt{type\_mandat}, \texttt{age\_moyen}, \texttt{age\_min}, \texttt{age\_max}, \texttt{q1\_age}, \texttt{mediane\_age}, \texttt{q3\_age}, \texttt{nb\_elus}

\begin{enumerate}[leftmargin=*]
  \item Glissez \texttt{type\_mandat} dans \textbf{Colonnes}.
  \item Glissez \texttt{age\_moyen} dans \textbf{Lignes}.
  \item Ajoutez \texttt{q1\_age}, \texttt{mediane\_age}, \texttt{q3\_age}, \texttt{age\_min}, \texttt{age\_max} dans \textbf{Info-bulle}.
  \item Pour simuler une boîte à moustaches : créez des \textbf{lignes de référence} (menu \textit{Analyse $\to$ Ligne de référence}) pour chaque statistique.
  \item Affichez \texttt{nb\_elus} dans \textbf{Taille} pour pondérer visuellement chaque type de mandat.
  \item Renommez la feuille \textit{``Distribution d'âge''}.
\end{enumerate}

\begin{hint}
Tableau Desktop ne génère pas de boîte à moustaches native à partir d'une table déjà agrégée (Q1/Q3 précalculés). Utilisez des \textbf{Lignes de référence} ou un graphique de type \textbf{Gantt} en encodant \texttt{q1\_age} sur l'axe Y et \texttt{q3\_age} en taille pour simuler les boîtes.
\end{hint}

\begin{validation}
Pour chaque type de mandat : l'âge moyen, le min, le max, Q1, la médiane et Q3 sont visibles. Les mandats sont comparables visuellement.
\end{validation}
\end{exercice}

% ── 5.3 Dashboard ────────────────────────────────────────────────────────────
\subsection*{5.3 — Assemblage du Dashboard}

\begin{exercice}{5e — Dashboard « Analyse des Élus RNE »}
\begin{enumerate}[leftmargin=*]
  \item Cliquez sur l'icône \textbf{Nouveau tableau de bord} (en bas de l'interface).
  \item Définissez la taille à \textbf{1200 $\times$ 800 px} (menu \textit{Taille du tableau de bord}).
  \item Glissez les 4 feuilles dans la disposition suivante :

\begin{center}
\begin{tabular}{|p{5.5cm}|p{5.5cm}|}
  \hline
  \textit{Parité H/F par mandat} & \textit{CSP par mandat} \\
  (barre empilée) & (treemap) \\
  \hline
  \textit{Ancienneté par région} & \textit{Distribution d'âge} \\
  (carte) & (boîte à moustaches) \\
  \hline
\end{tabular}
\end{center}

  \item Ajoutez un \textbf{titre} : \textit{``Répertoire National des Élus — Analyse des mandats''}.
  \item Rendez le filtre \texttt{Type Mandat} (de la feuille \textit{Ancienneté par région}) \textbf{global} : clic droit sur le filtre $\to$ \textbf{Appliquer aux feuilles de calcul $\to$ Toutes utilisant cette source de données}.
  \item Activez les \textbf{actions de filtre} : cliquer sur un type de mandat dans le graphique \textit{Parité} doit filtrer les autres vues. Menu \textbf{Tableau de bord $\to$ Actions $\to$ Ajouter une action $\to$ Filtre}.
  \item Renommez le tableau de bord \textit{``Dashboard RNE''}.
\end{enumerate}

\begin{validation}
\begin{itemize}[leftmargin=*]
  \item Les 4 vues sont présentes et lisibles dans le dashboard.
  \item Le filtre \texttt{Type Mandat} met à jour simultanément toutes les vues.
  \item Cliquer sur une barre du graphique \textit{Parité} filtre les autres visualisations.
\end{itemize}
\end{validation}
\end{exercice}

% =============================================================================
\section*{Récapitulatif}
% =============================================================================

\begin{tcolorbox}[colback=bleu!5, colframe=bleu, fonttitle=\bfseries,
                  title=Fichiers à rendre]
\begin{tabular}{ll}
  \texttt{spark\_jobs/job\_extract.py}   & Extraction PostgreSQL $\to$ HDFS \\
  \texttt{spark\_jobs/job\_transform.py} & Unification $\to$ \texttt{elus\_unified} \\
  \texttt{spark\_jobs/job\_aggregate.py} & 4 agrégations $\to$ HDFS + PostgreSQL \\
  \texttt{dags/rne\_pipeline.py}         & DAG Airflow avec 5 tâches \\
  \texttt{dashboard\_rne.twbx}           & Workbook Tableau (packaged) avec 4 vues + dashboard \\
\end{tabular}
\end{tcolorbox}

\vspace{0.8em}
\begin{tcolorbox}[colback=vertTP!5, colframe=vertTP, fonttitle=\bfseries,
                  title=Critères de validation]
\begin{enumerate}
  \item Les 5 tâches du DAG sont \textcolor{vertTP}{\textbf{vertes}} dans Airflow.
  \item Les 4 tables \texttt{agg\_*} sont présentes et non vides dans PostgreSQL.
  \item Les zones \texttt{/rne/raw}, \texttt{/rne/refined}, \texttt{/rne/gold} existent sur HDFS.
  \item Le fichier \texttt{.twbx} contient les 4 feuilles de visualisation et le dashboard assemblé.
  \item Le filtre \texttt{Type Mandat} dans le dashboard est global et interactif.
\end{enumerate}
\end{tcolorbox}

\end{document}
